% !PRINTED 479
\weberpage{479}
%<<CONT
oder in unserer abgekürzten Schreibweise:

\[
\left(\frac{\partial \Phi}{\partial x_1}, \frac{\partial \Phi}{\partial x_2}, \frac{\partial \Phi}{\partial x_3}\right) = A_1 \left(\frac{\partial \Phi}{\partial y_1}, \frac{\partial \Phi}{\partial y_2}, \frac{\partial \Phi}{\partial y_3}\right).
\]

Dies beweist den Satz:

1. \gesperrt{Die Variablenreihen}

\[
(x_1, x_2, x_3) \quad \text{und} \quad \left(\frac{\partial \Phi}{\partial x_1}, \frac{\partial \Phi}{\partial x_2}, \frac{\partial \Phi}{\partial x_3}\right)
\]

\gesperrt{sind contragredient.}

Durch wiederholte Anwendung dieses Satzes lassen sich auch die höheren Differentialquotienten nach den $x$ durch die nach den $y$ bilden, wofür man folgende Regel erhält:

2. \gesperrt{Um die $m^{\text{ten}}$ Ableitungen}

\[
\frac{\partial^m \Phi}{\partial x_1^\alpha \partial x_2^\beta \partial x_3^\gamma}, \quad \alpha + \beta + \gamma = m
\]

\gesperrt{durch die Ableitungen nach $y$ auszudrücken, ersetze man in dem entwickelten Ausdrucke}

\weq{4}{\xi_1^\alpha \xi_2^\beta \xi_3^\gamma = \sum^{\varkappa, \lambda, \mu} C^{\alpha, \beta, \gamma}_{\varkappa, \lambda, \mu} \eta_1^\varkappa \eta_2^\lambda \eta_3^\mu, \quad \varkappa + \lambda + \mu = m}

\gesperrt{die Producte}

\[
\begin{aligned}
&\xi_1^\alpha \xi_2^\beta \xi_3^\gamma \quad \text{durch} \quad \frac{\partial^m \Phi}{\partial x_1^\alpha \partial x_2^\beta \partial x_3^\gamma} \\
&\eta_1^\varkappa \eta_2^\lambda \eta_3^\mu \quad \text{durch} \quad \frac{\partial^m \Phi}{\partial y_1^\varkappa \partial y_2^\lambda \partial y_3^\mu},
\end{aligned}
\]

\gesperrt{also}

\weq{5}{\frac{\partial^m \Phi}{\partial x_1^\alpha \partial x_2^\beta \partial x_3^\gamma} = \sum C^{\alpha, \beta, \gamma}_{\varkappa, \lambda, \mu} \frac{\partial^m \Phi}{\partial y_1^\varkappa \partial y_2^\lambda \partial y_3^\mu},}

\gesperrt{wo unter dem Summenzeichen $\varkappa, \lambda, \mu$ alle nicht negativen der Bedingung $\varkappa + \lambda + \mu = m$ genügenden Werthe durchlaufen.}

Die Co\"efficienten $C^{\alpha, \beta, \gamma}_{\varkappa, \lambda, \mu}$ sind ganze rationale Functionen der Substitutionscoëfficienten $a_1, a_2, \ldots$

Um diese Regel allgemein zu beweisen, braucht man nur die Formel (5) für $m - 1$ statt $m$ als bewiesen anzusehen, und mit Anwendung der Formel (3) die Ableitung nach einer der Variablen $x$ zu bilden, und dabei die aus der Definition (4) folgende Relation

\[
C^{\alpha, \beta, \gamma}_{\varkappa, \lambda, \mu} = C^{\alpha - 1, \beta, \gamma}_{\varkappa - 1, \lambda, \mu}\, a_1 + C^{\alpha - 1, \beta, \gamma}_{\varkappa, \lambda - 1, \mu}\, a_2 + C^{\alpha - 1, \beta, \gamma}_{\varkappa, \lambda, \mu - 1}\, a_3
\]

zu berücksichtigen.
